\begin{lemma}
\label{lemma-hom-derived}
Let $(A, \text{d})$ be a differential graded algebra.
Let $M$ and $N$ be differential graded $A$-modules.
\begin{enumerate}
\item Let $P \to M$ be a P-resolution as in
Lemma \ref{lemma-resolve}. Then
$$
\Hom_{D(A, \text{d})}(M, N) =
\Hom_{K(\text{Mod}_{(A, \text{d})})}(P, N)
$$
\item Let $N \to I$ be an I-resolution as in
Lemma \ref{lemma-right-resolution}. Then
$$
\Hom_{D(A, \text{d})}(M, N) =
\Hom_{K(\text{Mod}_{(A, \text{d})})}(M, I)
$$
\end{enumerate}
\end{lemma}

\begin{proof}
Let $P \to M$ be as in (1). Since $P \to M$ is a quasi-isomorphism we see that
$$
\Hom_{D(A, \text{d})}(P, N) = \Hom_{D(A, \text{d})}(M, N)
$$
by definition of the derived category. A morphism
$f : P \to N$ in $D(A, \text{d})$ is equal to
$s^{-1}f'$ where $f' : P \to N'$ is a morphism and
$s : N \to N'$ is a quasi-isomorphism. Choose a distinguished triangle
$$
N \to N' \to Q \to N[1]
$$
As $s$ is a quasi-isomorphism, we see that $Q$ is acyclic. Thus
$\Hom_{K(\text{Mod}_{(A, \text{d})})}(P, Q[k]) = 0$ for all $k$ by
Lemma \ref{lemma-property-P-K-projective}. Since
$\Hom_{K(\text{Mod}_{(A, \text{d})})}(P, -)$
is cohomological, we conclude that we can lift $f' : P \to N'$
uniquely to a morphism $f : P \to N$. This finishes the proof.

\medskip\noindent
The proof of (2) is dual to that of (1) using
Lemma \ref{lemma-property-I-K-injective} in stead of
Lemma \ref{lemma-property-P-K-projective}.
\end{proof}
